% Generated draft: requires mathematical and source review.
\documentclass[11pt]{article}
\usepackage[margin=25mm]{geometry}
\usepackage{fontspec}
\setmainfont{Latin Modern Roman}
\usepackage{amsmath,amssymb,mathtools,unicode-math}
\setmathfont{Latin Modern Math}
\usepackage{enumitem,needspace,xurl,hyperref}
\setlength{\parindent}{0pt}
\setlength{\parskip}{5pt}
\setlength{\emergencystretch}{3em}
\newcommand{\N}{\mathbb{N}}
\newcommand{\Z}{\mathbb{Z}}
\newcommand{\Q}{\mathbb{Q}}
\newcommand{\R}{\mathbb{R}}
\newcommand{\C}{\mathbb{C}}
\newcommand{\F}{\mathbb{F}}
\newcommand{\A}{\mathbb{A}}
\newcommand{\PP}{\mathbb{P}}
\newcommand{\E}{\mathbb{E}}
\newcommand{\eps}{\varepsilon}
\newcommand{\dd}{\,\mathrm{d}}
\begin{document}
\section*{UnsolvedMath ID: 20001878}
\textbf{Source problem code:} catalogue entry
\subsection{Definitions and mathematical statement}
Hingston's Question 3.2.1 asks for a metric on \(S^n\) and an integral free-loop homology class \(X\) with
\[
 0<\operatorname{cr}(mX)<\operatorname{cr}(X)
 \quad\text{for some }m\in\N.
\]
Its simplest stated case is a generator \(X\in H_1(\Lambda S^2;\Z)\cong\Z\). Use \(\Lambda M=W^{1,2}(S^1,M)\), and fix the length-sublevel convention \(\Lambda^aM=\{\gamma:\operatorname{Length}(\gamma)\leq a\}\). Define
\[
 \operatorname{cr}(X)=\inf\{a:X\in\operatorname{im}[H_k(\Lambda^aM;\Z)\to H_k(\Lambda M;\Z)]\}.
\]
The source does not specify length versus energy normalization; the length choice is needed for the numerical lower bound below. Zero or annihilated torsion classes cannot satisfy the positive inequality.

The source relates this to the global mean frequency \(\alpha_{g,A}\), when it exists, defined by the limit of \(\deg Z/\operatorname{cr}_A(Z)\) over nonzero homogeneous classes as degree tends to infinity. It cites [Hin-Rad2013] for existence on spheres with field coefficients and asks whether the frequency or critical levels depend on the coefficient field. Its shorthand description of integral loop homology as \(\Z\) or \(\Z_2\) is not used here as a complete degreewise homology calculation; exceptional degrees or direct-sum issues require checking. The aid treats a primitive infinite-order class rather than resolving all coefficient behavior.

For any nested integral filtration, put \(I_a=\operatorname{im}H_k(\Lambda^aM;\Z)\). On the infinite cyclic subgroup generated by \(X\), there is a unique integer \(d(a)\geq0\) with
\[
 I_a\cap\Z X=d(a)\Z X,
\]
where \(d(a)=0\) means no nonzero multiple appears. The reported exact formulas are
\[
 \operatorname{cr}(mX)=\inf\{a:d(a)>0,\ d(a)\mid m\},
\]
\[
 \operatorname{cr}(\gcd(m,n)X)
 =\max\{\operatorname{cr}(mX),\operatorname{cr}(nX)\}.
\]
In particular a strict drop is equivalent to finding a level strictly below \(\operatorname{cr}(X)\) at which some nonzero multiple of \(X\), but not \(X\), lies in \(I_a\). Equivalently its image in the filtered quotient has finite nonzero order at such a separated level. Nonsaturation at a single endpoint alone need not force a strict difference of infima unless the filtration has an appropriate attainment/persistence property; that qualification is essential for an arbitrary nested filtration.

The aid also reports the rational envelope
\[
 \operatorname{cr}_{\Q}(X\otimes1)
 =\inf_{m\geq1}\operatorname{cr}_{\Z}(mX).
\]
It is a statement about this class under rational change of coefficients, not a proof that two field-valued global frequencies differ.

The source proposes a hypothetical four-level profile with \(a<b<c<d\): critical value \(d\) when \(m\) is coprime to both \(3,7\), \(c\) when \(3\mid m,7\nmid m\), \(b\) when \(7\mid m,3\nmid m\), and \(a\) when \(21\mid m\). The gcd law rules this out for one fixed class: \(\operatorname{cr}(X)=\max\{\operatorname{cr}(3X),\operatorname{cr}(7X)\}=c\), contradicting the proposed value \(d\). It was an illustrative pattern, not an asserted geometric example.

Finally, for a smooth Riemannian \(S^2\) the aid reports
\[
 \operatorname{cr}(mX)\geq2\operatorname{inj}(S^2,g)>0
 \quad(m\geq1)
\]
for the \(H_1\) generator in the length convention. Thus positivity is automatic in that case, but no metric producing the strict drop is constructed. The original existence and field-dependence questions remain unresolved in the supplied account.
\subsection{Short English statement}
Ask whether multiplying an integral loop-homology class can strictly lower its minimax level. The aid gives an exact divisor profile and gcd-maximum law, rules out the source's hypothetical independent 3/7 pattern, and proves positivity for the \(S^2\) generator. It supplies no metric realizing a strict drop.
\subsection{Sources}
\begin{itemize}
\item Local UnsolvedMath registry, permanent ID \(20001878\), code catalogue entry; canonical \texttt{aim-geometry-notes.json}, \texttt{notes[215]}; source tag ``question,'' statement status ``reconstructed\_unverified.''
\item Keith Burns and Vladimir S. Matveev, ``Open problems and questions about geodesics,'' Section3.2 ``Stability of minimax levels,'' Question3.2.1 attributed to Hingston, \url{https://aimath.org/pastworkshops/geodesicsproblems.pdf}. No PDF inspection is claimed here.
\item Source reference label [Hin-Rad2013], for the Resonance Theorem over fields, is retained without inventing its unavailable full bibliographic entry.
\item Research attempt 1: ``partial\_result,'' problem status at run ``open,'' registry status ``partially\_solved,'' rigor label ``proved,'' novelty confidence ``low.'' The research aid is machine generated, not peer reviewed; independently verify before citation.
\end{itemize}
\end{document}
